% =========================================================
% Approximation of Continuous Functions --- Notes
% Chapter: continuity
% Volume III - Analysis
% =========================================================

\subsection*{Approximation of Continuous Functions}
\label{sec:approximation}
% ---------------------------------------------------------
% Theorem --- Step Function Approximation
% ---------------------------------------------------------

\begin{theorembox}{Theorem (Step Function Approximation)}
\begin{theorem}[Step Function Approximation]
\label{thm:step-function-approximation}
Let $f:[a,b]\to\mathbb{R}$ be continuous and let $\varepsilon>0$.
Then there exists a step function $s_\varepsilon:[a,b]\to\mathbb{R}$ such that
\[
\sup_{x\in[a,b]} |f(x)-s_\varepsilon(x)|<\varepsilon.
\]
\hyperref[prf:step-function-approximation]{Go to proof.}
\end{theorem}
\end{theorembox}

\begin{remark*}[Standard quantified statement]
\[
\begin{aligned}
&\text{Let } f:[a,b]\to\mathbb{R} \text{ be continuous and let } \varepsilon>0.\\
&\exists s_\varepsilon:[a,b]\to\mathbb{R} \text{ step function } \forall x\in[a,b]\;(|f(x)-s_\varepsilon(x)|<\varepsilon).
\end{aligned}
\]
\end{remark*}

\begin{remark*}[Predicate reading]
\[
\operatorname{IsStepFunction}(s_\varepsilon,[a,b])
\land
\forall x\in[a,b]\;(|f(x)-s_\varepsilon(x)|<\varepsilon).
\]
\end{remark*}

\begin{remark*}[Negated quantified statement]
\[
\begin{aligned}
&\text{Let } f:[a,b]\to\mathbb{R} \text{ be continuous and let } \varepsilon>0.\\
&\forall s:[a,b]\to\mathbb{R} \text{ step function } \exists x\in[a,b]\;(|f(x)-s(x)|\ge\varepsilon).
\end{aligned}
\]
\end{remark*}

\begin{remark*}[Negation predicate reading]
\[
  s\colon A\to\mathbb{R},\\

\forall s:[a,b]\to\mathbb{R}\;
\Bigl(\operatorname{IsStepFunction}(s,[a,b]) \Rightarrow
\exists x\in[a,b]\;(|f(x)-s(x)|\ge\varepsilon)\Bigr).
\]
\end{remark*}

\begin{remark*}[Failure modes]
\begin{description}
  \item[Exposition.]
The statement fails precisely when a continuous $f$ and tolerance $\varepsilon$ are fixed but every step function leaves an error of at least $\varepsilon$ at some point in $[a,b]$.


  \item[\operatorname{IsStepFunction}.]
  The displayed obstruction is the corresponding failure mode.
  \[
\begin{aligned}
&\text{Let } f:[a,b]\to\mathbb{R} \text{ be continuous and let } \varepsilon>0.\\
&\forall s:[a,b]\to\mathbb{R} \text{ step function } \exists x\in[a,b]\;(|f(x)-s(x)|\ge\varepsilon).
\end{aligned}
  \]
  \[
\forall s:[a,b]\to\mathbb{R}\;
\Bigl(\operatorname{IsStepFunction}(s,[a,b]) \Rightarrow
\exists x\in[a,b]\;(|f(x)-s(x)|\ge\varepsilon)\Bigr).
  \]
\end{description}
\end{remark*}

\begin{remark*}[Failure modes]
\begin{description}
  \item[Exposition.]
\[
\underbrace{\forall s:[a,b]\to\mathbb{R} \text{ step function }}_{\text{all approximants}}
\underbrace{\exists x\in[a,b]\;(|f(x)-s(x)|\ge\varepsilon)}_{\text{witnessed large deviation}}.
\]


  \item[\operatorname{IsStepFunction}.]
  The displayed obstruction is the corresponding failure mode.
  \[
\begin{aligned}
&\text{Let } f:[a,b]\to\mathbb{R} \text{ be continuous and let } \varepsilon>0.\\
&\forall s:[a,b]\to\mathbb{R} \text{ step function } \exists x\in[a,b]\;(|f(x)-s(x)|\ge\varepsilon).
\end{aligned}
  \]
  \[
\forall s:[a,b]\to\mathbb{R}\;
\Bigl(\operatorname{IsStepFunction}(s,[a,b]) \Rightarrow
\exists x\in[a,b]\;(|f(x)-s(x)|\ge\varepsilon)\Bigr).
  \]
\end{description}
\end{remark*}

\begin{remark*}[Interpretation]
Uniform continuity on $[a,b]$ guarantees that $f$ does not oscillate rapidly at small scales, so subdividing the interval finely enough allows a step function to track $f$ within any prescribed uniform error. The typical construction samples $f$ at one point of each subinterval and keeps that value constant across the piece. Failure would mean that no matter how the interval is partitioned, some subinterval forces an error at least $\varepsilon$, contradicting uniform continuity. The result is a foundational approximation fact used to show that continuous functions on compact intervals are Riemann integrable and to connect continuous functions with simple combinatorial models.
\end{remark*}

\begin{dependencies}
\hyperref[thm:heine-cantor]{Heine--Cantor Theorem}.
\end{dependencies}

% ---------------------------------------------------------
% Definition --- Piecewise Linear Function
% ---------------------------------------------------------

\begin{definitionbox}{Definition (Piecewise Linear Function)}
\begin{definition}[Piecewise Linear Function]
\label{def:piecewise-linear-function}
Let $a,b \in \mathbb{R}$ with $a<b$ and let $g:[a,b]\to\mathbb{R}$.
The function $g$ is piecewise linear if and only if there exist
$m \in \mathbb{N}$ with $m \ge 1$ and points
$a=x_0<x_1<\cdots<x_m=b$ such that $g|_{[x_{k-1},x_k]}$ is affine
for every $k\in\{1,\ldots,m\}$.
\end{definition}
\end{definitionbox}

\begin{remark*}[Standard quantified statement]
\[
\begin{aligned}
&\text{Let } a,b \in \mathbb{R} \text{ with } a<b \text{ and } g:[a,b]\to\mathbb{R}.\\
&g \text{ is piecewise linear}\\
&\Longleftrightarrow
\exists m \in \mathbb{N}\;\Bigl(m \ge 1 \land
\exists x_0,\ldots,x_m \in [a,b]\;
\bigl[a=x_0 < x_1 < \cdots < x_m = b\\
&\qquad\qquad\land
\forall k \in \{1,\ldots,m\}\; g|_{[x_{k-1},x_k]} \text{ is affine}\bigr]\Bigr).
\end{aligned}
\]
\end{remark*}

\begin{remark*}[Predicate reading]
\[
  g\colon A\to\mathbb{R},\\

\operatorname{IsPiecewiseLinear}(g,[a,b]) \coloneqq
\exists m \in \mathbb{N}\;\Bigl(m \ge 1 \land
\exists x_0,\ldots,x_m \in [a,b]\;
\bigl[a=x_0 < \cdots < x_m = b \land
\forall k \in \{1,\ldots,m\}\; g|_{[x_{k-1},x_k]} \text{ is affine}\bigr]\Bigr).
\]
\end{remark*}

\begin{remark*}[Negated quantified statement]
\[
\begin{aligned}
&\text{Let } a,b \in \mathbb{R} \text{ with } a<b \text{ and } g:[a,b]\to\mathbb{R}.\\
&g \text{ is not piecewise linear}\\
&\Longleftrightarrow
\forall m \in \mathbb{N}\;\Bigl(m \ge 1 \Rightarrow
\forall x_0,\ldots,x_m \in [a,b]\;
\bigl[a=x_0 < \cdots < x_m = b\\
&\qquad\qquad\Rightarrow
\exists k \in \{1,\ldots,m\}\; g|_{[x_{k-1},x_k]} \text{ is not affine}\bigr]\Bigr).
\end{aligned}
\]
\end{remark*}

\begin{remark*}[Negation predicate reading]
\[
  g\colon A\to\mathbb{R},\\

\neg\operatorname{IsPiecewiseLinear}(g,[a,b])
\Longleftrightarrow
\forall m \in \mathbb{N}\;\Bigl(m \ge 1 \Rightarrow
\forall x_0,\ldots,x_m \in [a,b]\;
\bigl[a=x_0 < \cdots < x_m = b
\Rightarrow
\exists k \in \{1,\ldots,m\}\; g|_{[x_{k-1},x_k]} \text{ is not affine}\bigr]\Bigr).
\]
\end{remark*}

\begin{remark*}[Failure modes]
\begin{description}
  \item[Exposition.]
A proposed piecewise-linear structure can fail because no finite ordered
partition of $[a,b]$ is available, or because for every such partition
at least one subinterval has a restriction of $g$ that is not affine.


  \item[\operatorname{IsPiecewiseLinear}.]
  The displayed obstruction is the corresponding failure mode.
  \[
\begin{aligned}
&\text{Let } a,b \in \mathbb{R} \text{ with } a<b \text{ and } g:[a,b]\to\mathbb{R}.\\
&g \text{ is not piecewise linear}\\
&\Longleftrightarrow
\forall m \in \mathbb{N}\;\Bigl(m \ge 1 \Rightarrow
\forall x_0,\ldots,x_m \in [a,b]\;
\bigl[a=x_0 < \cdots < x_m = b\\
&\qquad\qquad\Rightarrow
\exists k \in \{1,\ldots,m\}\; g|_{[x_{k-1},x_k]} \text{ is not affine}\bigr]\Bigr).
\end{aligned}
  \]
  \[
\neg\operatorname{IsPiecewiseLinear}(g,[a,b])
\Longleftrightarrow
\forall m \in \mathbb{N}\;\Bigl(m \ge 1 \Rightarrow
\forall x_0,\ldots,x_m \in [a,b]\;
\bigl[a=x_0 < \cdots < x_m = b
\Rightarrow
\exists k \in \{1,\ldots,m\}\; g|_{[x_{k-1},x_k]} \text{ is not affine}\bigr]\Bigr).
  \]
\end{description}
\end{remark*}

\begin{remark*}[Failure modes]
\begin{description}
  \item[Exposition.]
\[
\underbrace{\forall m\in\mathbb{N}\;(m\ge 1\Rightarrow
\forall x_0,\ldots,x_m\in[a,b]\;(a=x_0<\cdots<x_m=b\Rightarrow\cdots))}_{\text{no valid finite affine subdivision}}
\;\lor\;
\underbrace{\exists k\in\{1,\ldots,m\}\; g|_{[x_{k-1},x_k]}\text{ is not affine}}_{\text{non-affine segment}}.
\]


  \item[\operatorname{IsPiecewiseLinear}.]
  The displayed obstruction is the corresponding failure mode.
  \[
\begin{aligned}
&\text{Let } a,b \in \mathbb{R} \text{ with } a<b \text{ and } g:[a,b]\to\mathbb{R}.\\
&g \text{ is not piecewise linear}\\
&\Longleftrightarrow
\forall m \in \mathbb{N}\;\Bigl(m \ge 1 \Rightarrow
\forall x_0,\ldots,x_m \in [a,b]\;
\bigl[a=x_0 < \cdots < x_m = b\\
&\qquad\qquad\Rightarrow
\exists k \in \{1,\ldots,m\}\; g|_{[x_{k-1},x_k]} \text{ is not affine}\bigr]\Bigr).
\end{aligned}
  \]
  \[
\neg\operatorname{IsPiecewiseLinear}(g,[a,b])
\Longleftrightarrow
\forall m \in \mathbb{N}\;\Bigl(m \ge 1 \Rightarrow
\forall x_0,\ldots,x_m \in [a,b]\;
\bigl[a=x_0 < \cdots < x_m = b
\Rightarrow
\exists k \in \{1,\ldots,m\}\; g|_{[x_{k-1},x_k]} \text{ is not affine}\bigr]\Bigr).
  \]
\end{description}
\end{remark*}

\begin{remark*}[Interpretation]
A piecewise linear function on $[a,b]$ has a graph made from finitely
many straight line segments. The definition permits different affine
rules on different subintervals, but requires a finite ordered list of
nodes that covers the whole interval.
\end{remark*}

\begin{dependencies}
\begin{itemize}
  \item \hyperref[def:natural-numbers]{Natural Numbers}
  \item \hyperref[def:order-on-r]{Order on $\mathbb{R}$}
  \item \hyperref[def:closed-interval]{Closed Interval}
  \item \hyperref[def:function]{Function}
\end{itemize}
\end{dependencies}
% ---------------------------------------------------------
% Theorem --- Piecewise Linear Approximation
% ---------------------------------------------------------

\begin{theorembox}{Theorem}
\begin{theorem}[Piecewise Linear Approximation]
\label{thm:piecewise-linear-approximation}
Let $f\colon[a,b]\to\mathbb{R}$ be continuous and let $\varepsilon>0$. Then there exists a continuous piecewise linear function $\ell_\varepsilon\colon[a,b]\to\mathbb{R}$ such that $\sup_{x\in[a,b]}|f(x)-\ell_\varepsilon(x)|<\varepsilon$. \hyperref[prf:piecewise-linear-approximation]{Go to proof.}
\end{theorem}
\end{theorembox}

\begin{remark*}[Standard quantified statement]
\[
\begin{aligned}
&\text{Let } f\colon[a,b]\to\mathbb{R},\; f \text{ continuous},\; \varepsilon>0.\\
&\exists\text{ continuous piecewise linear }\ell_\varepsilon\colon[a,b]\to\mathbb{R}:
\sup_{x\in[a,b]}|f(x)-\ell_\varepsilon(x)|<\varepsilon.
\end{aligned}
\]
\end{remark*}

\begin{remark*}[Predicate reading]
\[
\begin{aligned}
&\text{Let } f\colon[a,b]\to\mathbb{R},\; f \text{ continuous},\; \varepsilon>0.\\
&\exists\text{ continuous piecewise linear }\ell_\varepsilon\colon[a,b]\to\mathbb{R}:
\sup_{x\in[a,b]}|f(x)-\ell_\varepsilon(x)|<\varepsilon.
\end{aligned}
\]
\end{remark*}

\begin{remark*}[Interpretation]
Every continuous function on a compact interval can be uniformly approximated by a polygonal curve matching finitely many sampled values of the original graph. The construction refines the interval into a sufficiently fine partition guaranteed by compactness, then connects the sampled values with straight segments, producing a uniformly small error over the entire domain. Failure would require a uniform oscillation exceeding $\varepsilon$ on some subinterval, which cannot occur once the partition is fine enough. Geometrically, the theorem states that continuous graphs admit arbitrarily close polyline shadows, providing a bridge from rough step approximations to smooth polynomial approximations.
\end{remark*}

\begin{dependencies}
\begin{itemize}
  \item \hyperref[def:piecewise-linear-function]{Piecewise Linear Function}
  \item \hyperref[thm:heine-cantor]{Heine--Cantor Theorem}
  \item \hyperref[def:continuous-at-point]{Continuity at a Point}
\end{itemize}
\end{dependencies}
% ---------------------------------------------------------
% Theorem --- Weierstrass Approximation Theorem
% ---------------------------------------------------------

\begin{theorembox}{Theorem (Weierstrass Approximation Theorem)}
\begin{theorem}[Weierstrass Approximation Theorem]
\label{thm:weierstrass-approximation}
Let $f:[a,b]\to\mathbb{R}$ be continuous. For every $\varepsilon>0$ there exists a polynomial $p_\varepsilon\in\mathbb{R}[x]$ such that $\|f-p_\varepsilon\|_{\infty,[a,b]}<\varepsilon$.
\hyperref[prf:weierstrass-approximation]{Go to proof.}
\end{theorem}
\end{theorembox}

\begin{remark*}[Standard quantified statement]
\[
\begin{aligned}
&\text{Let } a<b \text{ in } \mathbb{R}, \text{ and let } f:[a,b]\to\mathbb{R} \text{ be continuous}.\\
&\forall \varepsilon>0\;\exists p\in\mathbb{R}[x]\;(\|f-p\|_{\infty,[a,b]}<\varepsilon).
\end{aligned}
\]
\end{remark*}

\begin{remark*}[Predicate reading]
\[
  f\colon A\to\mathbb{R},\\

\operatorname{IsContinuous}(f,[a,b],\mathbb{R},\mathbb{R})
\Longrightarrow
\forall \varepsilon>0\;\exists p\in\mathbb{R}[x]\;
\bigl(\|f-p\|_{\infty,[a,b]}<\varepsilon\bigr).
\]
\end{remark*}

\begin{remark*}[Interpretation]
Every continuous function on a closed interval can be uniformly approximated by polynomials, so polynomials are dense in the space of continuous functions under the supremum norm. The result holds because convolutions with Bernstein polynomials preserve continuity and converge uniformly to the original function. Failure would mean there exists a uniform tolerance that no polynomial can match, contradicting the approximation scheme built from Bernstein polynomials. Geometrically, the theorem states that the graph of a continuous function can be matched arbitrarily well by the graph of a polynomial, so polynomial graphs form a flexible mesh that can shadow any continuous curve on a compact interval.
\end{remark*}

\begin{dependencies}
Depends on \hyperref[thm:bernstein-approximation]{thm:bernstein-approximation}.
\end{dependencies}

% ---------------------------------------------------------
% Definition --- Bernstein Polynomial
% ---------------------------------------------------------

\begin{definitionbox}{Definition (Bernstein Polynomial)}
\begin{definition}[Bernstein Polynomial]
\label{def:bernstein-polynomial}
Let $n \in \mathbb{N}$ and let $f:[0,1]\to\mathbb{R}$.
A function $B_n:[0,1]\to\mathbb{R}$ is called the $n$th Bernstein polynomial of $f$
if and only if, for every $x \in [0,1]$,
\[
B_n(x)=\sum_{k=0}^{n} f\!\left(\frac{k}{n}\right)\binom{n}{k}x^k(1-x)^{n-k}.
\]
\end{definition}
\end{definitionbox}

\begin{remark*}[Standard quantified statement]
\[
\begin{aligned}
&\text{Let }n\in\mathbb{N},\; f:[0,1]\to\mathbb{R},\; B_n:[0,1]\to\mathbb{R}.\\
&B_n \text{ is the } n\text{th Bernstein polynomial of } f\\
&\Longleftrightarrow
\forall x\in[0,1]\;
\Bigl(
B_n(x)=\sum_{k=0}^n f(k/n)\binom{n}{k}x^k(1-x)^{n-k}
\Bigr).
\end{aligned}
\]
\end{remark*}

\begin{remark*}[Predicate reading]
\[
\Bigl(B_n:[0,1]\to\mathbb{R}\Bigr)\land
\forall x\in[0,1]\;
\Bigl(
B_n(x)=\sum_{k=0}^n f(k/n)\binom{n}{k}x^k(1-x)^{n-k}
\Bigr).
\]
\end{remark*}

\begin{remark*}[Interpretation]
The nth Bernstein polynomial samples f at the uniformly spaced nodes k/n, weights those samples by the binomial probabilities of order n, and produces for each x an average that depends polynomially on x. This construction smooths f into a polynomial that agrees with f at the endpoints and approximates f uniformly as n grows, so describing it explicitly is essential for subsequent approximation theorems. Failure arises only when the candidate polynomial disagrees with the binomial averaging recipe at some point of the unit interval, reflecting an incorrect coefficient or evaluation in the polynomial expression.
\end{remark*}

\begin{dependencies}
\begin{itemize}
  \item \hyperref[def:natural-numbers]{Natural Numbers}
  \item \hyperref[def:addition-on-r]{Addition on $\mathbb{R}$}
  \item \hyperref[def:function]{Function}
  \item \hyperref[def:multiplication-on-r]{Multiplication on $\mathbb{R}$}
\end{itemize}
\end{dependencies}
% ---------------------------------------------------------
% Theorem --- Bernstein Approximation Theorem
% ---------------------------------------------------------

\begin{theorembox}{Theorem (Bernstein Approximation Theorem)}
\begin{theorem}[Bernstein Approximation Theorem]
\label{thm:bernstein-approximation}
Let $f:[0,1]\to\mathbb{R}$ be continuous, and for each $n\in\mathbb{N}$ let $B_n$ denote the $n$th Bernstein polynomial of $f$. Then for every $\varepsilon>0$ there exists $n_\varepsilon\in\mathbb{N}$ such that $\|f-B_n\|_\infty<\varepsilon$ for all $n\ge n_\varepsilon$. \hyperref[prf:bernstein-approximation]{Go to proof.}
\end{theorem}
\end{theorembox}

\begin{remark*}[Standard quantified statement]
\[
\begin{aligned}
&\text{Let } f:[0,1]\to\mathbb{R} \text{ be continuous, and for each } n\in\mathbb{N} \text{ let } B_n \text{ be the $n$th Bernstein polynomial of } f.\\
&\forall \varepsilon>0\;\exists n_\varepsilon\in\mathbb{N}\;\forall n\ge n_\varepsilon:\ \|f-B_n\|_\infty<\varepsilon.
\end{aligned}
\]
\end{remark*}

\begin{remark*}[Predicate reading]
\[
\begin{aligned}
&\text{Let } f:[0,1]\to\mathbb{R} \text{ be continuous, and for each } n\in\mathbb{N} \text{ let } B_n \text{ be the $n$th Bernstein polynomial of } f.\\
&\forall \varepsilon>0\;\exists n_\varepsilon\in\mathbb{N}\;\forall n\ge n_\varepsilon:\ \|f-B_n\|_\infty<\varepsilon.
\end{aligned}
\]
\end{remark*}

\begin{remark*}[Interpretation]
The theorem asserts that the Bernstein polynomials of a continuous function on $[0,1]$ converge uniformly to the original function, so the entire graph of $f$ can be approximated to any prescribed sup-norm accuracy by explicit degree-$n$ polynomials once $n$ is large enough. The proof works by showing that the Bernstein operators preserve positivity and linear growth while averaging values of $f$, which forces the uniform error to vanish as $n\to\infty$. This matters because it gives an elementary constructive route from continuous functions to polynomials, strengthening the Weierstrass approximation theorem with a concrete approximating sequence. The approximation can fail only when $f$ is not continuous on $[0,1]$, in which case the Bernstein polynomials may smooth out jumps and no longer capture the original function in the sup norm. Geometrically, each $B_n$ is a convex combination of the samples $f(k/n)$, so the polygons formed by these samples become finer and eventually hug the graph of $f$ everywhere on the interval.
\end{remark*}

\begin{dependencies}
\hyperref[def:continuous-at-point]{Definition (Continuity at a Point)}, \hyperref[def:bernstein-polynomial]{Definition (Bernstein Polynomial)}.
\end{dependencies}
